\section{Effective nerve theorems for finite covers}
\label{sec:nerve-case}

The repository's topology report lists an effective nerve theorem with
produced contraction data among its open tasks. Its carrier calculus handles
ambiguity of maps between covers. Even for one cover we must determine when its incidence pattern recovers the chains in its union. We give two algorithms. The first
tests the exact finite acyclicity hypothesis and constructs chain equivalences
over $\mathbb F_2$. The second produces the local data directly from the
endpoints of rational boxes. Both give explicit chain maps in both directions.
The classical nerve theorems for closed convex covers and subcomplex covers
are treated by Bauer, Kerber, Roll and Rolle. Here the additional arithmetic
content is the production of bounded filling tables, homotopies and
moduli. See the
\repo{docs/proofs/topology/ipra_topology_framework_2026-11-15/paper/sections/12_limits.tex}{repository's topology problem list}
and \href{https://arxiv.org/abs/2203.03571}{Bauer et al.,\ 2023}.

\subsection{The input and the local computation}

Let $K$ be a nonempty finite simplicial or elementary grid-cubical complex,
given by cell and face tables, covered by a finite list of nonempty
subcomplexes $U_1,\ldots,U_m$. Cover means that every cell occurs in some
$U_i$. For nonempty $I\subseteq\{1,\ldots,m\}$ put
$K_I=\bigcap_{i\in I} U_i$. The nerve $L$ consists of the $I$ for which
$K_I$ contains a cell. We use ordinary chains unless they are expressly
augmented. An $r$-chain is a bit vector indexed by the $r$-cells. Add by XOR
and extend each boundary column linearly. For a vertex its ordinary boundary
is zero. The augmented complex has one additional basis element $*$ in
degree $-1$, and $\partial v=*$ at every vertex. The identity
$\partial^2=0$ follows by pairing the two orders in which two
vertices are removed or two distinct coordinate faces are taken. In degree
one the two endpoints have the same augmentation.

\paragraph{Local filling.}
In each nonempty $K_I$ assign $h(*)$ its first vertex. Process cells $c$ in
increasing dimension. Set $z=c+h\partial c$ and seek a vector $u$ of the
next degree satisfying $\partial u=z$. On success set $h(c)=u$. The
previous equations give $\partial h\partial c=\partial c$,
so $\partial z=0$. If every augmented cycle bounds, the finite search
therefore succeeds. Conversely, success on all cells gives
$\partial h+h\partial=1$ by columns, hence any augmented cycle $z'$ bounds
$h(z')$. On failure the algorithm has produced the obstructing cycle $z$ after testing every possible filling. A group of $t$ possible basis cells
requires at most $2^t$ trials. This proves the equivalence between the test
and vanishing of all reduced homology of the intersection, including
connectedness in degree zero.

\paragraph{Finite nerve theorem.}
If the local test succeeds for every simplex of $L$, a primitive recursive
algorithm constructs a chain homotopy equivalence
$C_*(K;\mathbb F_2)\simeq C_*(L;\mathbb F_2)$ together with the two
homotopies witnessing its inverse laws. The test uses homological
acyclicity. The conclusion concerns chains and homology.

\subsection{Eliminating the intersections}

Form an incidence complex $E$ with basis $(I,c)$, where $c$ is any cell of
$K_I$. Its bidegree is $(p,q)=(|I|-1,\dim c)$ and its ordinary total
degree is $p+q$. Its horizontal boundary $b$ deletes one index from $I$. Its vertical boundary $v$ takes the boundary of $c$. Each is zero at its
ordinary degree zero. They commute because every face of $c$ belongs to
every subcomplex containing $c$. Thus $D=b+v$ squares to zero in
characteristic two. Define the projections by
\[
 \pi_K(I,c)=
 \begin{cases}c& |I|=1,\\0&\text{otherwise,}\end{cases}
 \qquad
 \pi_L(I,c)=
 \begin{cases}I& \dim c=0,\\0&\text{otherwise.}\end{cases}
\]
They commute with the respective differentials. At the first nonzero
horizontal degree the two deletions project to the same cell twice. In
vertical degree one the two endpoints project to the same index simplex twice.

For $\pi_K$ append a column $p=-1$ having basis $(\varnothing,c)$, and
allow $b$ to delete the last index. Give the augmented array degree $p+q+1$,
so the appended column is $C_*(K)$ in its ordinary degree. Fix $c$ and let
$j$ be the least chart containing it. Its horizontal row is the augmented
simplex on the nonempty set of charts containing $c$. Define $h(I,c)$ by
adjoining $j$ when absent and by zero otherwise. If $j\notin I$, all
deletions other than deletion of $j$ cancel against $hb$. If $j\in I$,
only its deletion contributes in $hb$. This includes $I=\varnothing$
and proves $bh+hb=1$.

Consequently
\[
 R=1+Dh+hD=vh+hv.
\]
Each nonzero summand lowers $q$ by one, giving $R^{e+1}=0$ for
$e=\dim K$. Also $DR=RD$, by expansion and $D^2=0$. With
$H=h(1+R+\cdots+R^e)$ we obtain
\[
 DH+HD=(1+R)(1+R+\cdots+R^e)=1+R^{e+1}=1.
\]
For $\pi_L$, append instead the row $q=-1$ on $(I,*)$ and augment the
vertical differential. Take the previously constructed $h$ of $K_I$ on each
column. Now $vh+hv=1$ and each term of $R=bh+hb$ lowers $p$. Use the sum of
powers through $\dim L$. Both augmented arrays thus have explicit global
contractions, independently of the choices of local fillings.

To read the result, write either augmented array in degree $n$ as
$B_n\oplus E_{n-1}$, with $B=C_*(K)$ or $C_*(L)$. Its differential and
contracting homotopy have block form
\[
 \begin{pmatrix}d_B&\pi\\0&d_E\end{pmatrix},
 \qquad
 \begin{pmatrix}a&b'\\s&t\end{pmatrix}.
\]
Three blocks of $DH+HD=1$ say that $s:B\to E$ is a chain map,
$1+\pi s=d_Ba+ad_B$, and $1+s\pi=d_Et+td_E$. This supplies both
inverse laws. Denote the two sets of blocks by subscripts $K,L$.
The maps $f=\pi_L s_K$ and $g=\pi_K s_L$ go between the requested
complexes. A homotopy for $1+gf$ is
$a_K+\pi_Kt_Ls_K$. The other is $a_L+\pi_Lt_Ks_L$.
Substitution proves the formulas even when $K$ is disconnected or the nerve
contains repeated charts as distinct vertices.

\subsection{Rectangular sets of any dimension}

Consider a finite nonempty list of closed boxes
$B_i=\prod_{k=1}^d[a_{ik},b_{ik}]$ with rational endpoints
$a_{ik}\le b_{ik}$. Sort and deduplicate all endpoints separately in each
coordinate. The products of grid vertices and intervening elementary edges
that lie in some box, together with their faces, form $K$. Their closed
geometric cells cover exactly the union of the boxes, since each coordinate in a
closed box lies at a listed endpoint or in an intervening interval whose
closed endpoints still belong to that box. This is a geometric reading
for classical points. The internal input consists entirely of rational
endpoints and finite cells. The intersection indexed by
$I$ is present exactly when, in every coordinate,
$\max_{i\in I}a_{ik}\le\min_{i\in I}b_{ik}$. Equality permits a
lower-dimensional overlap. This formula computes the nerve from the original
list.

Here the vertical contraction has an explicit formula. In a one-dimensional
intersection with grid vertices $w_0,\ldots,w_l$ in order, let $P$
send each vertex to $w_0$ and each edge to zero. Let $u(w_j)$ be the sum of
the first $j$ edges and let $u$ vanish on edges. Consecutive endpoints
cancel, giving $\partial u+u\partial=1+P$ in the ordinary complex,
also for $l=0$. In the product grid a cell is a tensor of vertices and
edges. Put
\[
 u_{\mathrm{box}}=
 \sum_{k=1}^d P_1\otimes\cdots\otimes P_{k-1}\otimes
 u_k\otimes 1\otimes\cdots\otimes 1.
\]
The product differential acts in one coordinate at a time. Terms acting in a
coordinate different from $k$ cancel in pairs since each $P_i$ commutes
with boundary. The remaining term for $k$ is
$(P_1\otimes\cdots\otimes P_{k-1})\otimes(1+P_k)\otimes
 1\otimes\cdots\otimes 1$. Successive summands cancel their common
partial projections, leaving $1+P_1\otimes\cdots\otimes P_d$.
The final projection kills all positive-dimensional cells and sends every
vertex to the lower corner. Sending $*$ to this corner extends the formula
to a contraction of the augmented complex. For $d=0$ the single corner
itself furnishes the augmentation inverse. This finishes the construction
from endpoints alone. The topological invariant calculated internally is
grid-cubical homology. Identifying it with singular homology of the realized
polyhedron is the external cellular comparison.

\subsection{The full homotopy type of a union of rational boxes}

For these boxes there is a stronger conclusion. Give $X=\bigcup_i B_i$
its geometric grid realization. Realize $L$ by probability vectors $y$
whose positive support is a simplex of $L$. There are explicit rational
piecewise-affine maps $F:X\to |L|$ and $G:|L|\to X$ and effective
homotopies from both composites to the respective identities. We construct
them on rational points, bound their moduli, and then read them on regular
names. The geometric argument avoids passing through an abstract nerve
theorem.

In the barycentric subdivisions of $K$ and $L$, vertices are
barycenters of nonempty cells, and a simplex is a strict chain of such cells.
For a box cell or simplex, coning its already subdivided boundary to its
center gives this triangulation by dimension induction. Each point distinct
from the center is on a unique segment from the center to a boundary point,
as can be checked along the ray by the first bounding face it meets. The
triangulations agree on common faces. For each grid cell $c$, let $I(c)$
be its nonempty set of covering labels. Write $e_I$ for the probability
vector uniform on $I$. Set
\[
    F(\operatorname{bar}c)=e_{I(c)},\qquad
    G(e_I)=\ell_I,\qquad
    (\ell_I)_k=\max_{i\in I}a_{ik}.
\]
Extend affinely on each simplex of the respective subdivisions. Along a
cell chain, the sets $I(c)$ are nested in reverse order, so $F$ takes
values in $|L|$. Along an index chain $I_0\subset\cdots\subset I_j$,
the corners $\ell_{I_t}$ all belong to $B_{I_0}=\bigcap_{i\in I_0}B_i$.
Their convex combinations belong there too. Common-face agreement follows
from affine interpolation.

If $x$ lies on a simplex with largest cell $c$, then $x$ and $GF(x)$
both lie in $B_{I(c)}$. Indeed every $I(c')$ occurring in its image
contains $I(c)$, and its lower corner belongs to each box in $I(c)$.
Thus the straight segment from $x$ to $GF(x)$ stays in $X$.

The other composite needs more care since the nerve realization need not be
convex. Write $z=FG(y)$. In a subdivision simplex representing $y$, discard
the zero barycentric coefficients and call the smallest surviving label set
$I_0$. For any $i\in I_0$, every vertex contributing to $y$ gives label $i$
weight at least $1/m$, hence $y_i\ge 1/m$. Also $G(y)\in B_i$. For
\emph{every} point $x\in B_i$, the inequality $F_i(x)\ge1/m$ holds because the grid cells
contained in $B_i$ form a subcomplex covering it, and the barycentric
interpolation there uses only vertices whose covering label sets contain $i$.
This proves $z_i\ge1/m$ even on shared faces, where multiple simplex
representations are available. Define
\[
 W_i(y)=\frac{y_i z_i}{\sum_{j=1}^{m}y_jz_j}.
\]
The denominator is at least $1/m^2$. The support of $W$ is contained in
each of the two supports, so the segment from $y$ to $W$ lies in the
simplex supporting $y$, and the segment from $W$ to $z$ lies in the simplex
supporting $z$. Traversing each segment in half the time gives a continuous
homotopy. It agrees across all subdivision boundaries because the displayed
formula depends only on $y$ and $z$. This establishes the full homotopy
equivalence for this geometric class.

\paragraph{Evaluation and moduli.}
The internal calculations use rational points. A simplex of the barycentric
triangulation containing a given rational point can be selected by bounded
search over all chains and exact barycentric tests. More directly, for a grid
cell $c$ with center $b$, the fraction along the center-to-boundary ray is
$\max_k 2|x_k-b_k|/\operatorname{width}_k(c)$ over nonconstant coordinates.
At fraction zero evaluate at the center. At a positive fraction interpolate and recurse
on a proper face. To evaluate $G$, list the distinct positive entries
$t_1>\cdots>t_r$ of $y$ and set $t_{r+1}=0$. Use the convex coefficients
$|I_j|(t_j-t_{j+1})$ at the corners $\ell_{I_j}$ with
$I_j=\{i:y_i\ge t_j\}$. They sum to one because their contributions to
coordinate $i$ in the original convex combination of the $e_{I_j}$
telescope to $y_i$.

Here is a uniform PR recipe for Lipschitz bounds in the sup norm. Extend
$F$ by assigning zero to the unused cell barycenters in the entire
bounding grid and interpolating there. Extend $G$ to the entire probability
simplex in the same way, assigning a fixed grid point to missing index
simplices. These are continuous PL extensions to convex polytopes. On each
maximal simplex compute the rational affine coefficient matrix by solving
the vertex equations. Omit constant grid coordinates and use the first
$m-1$ independent probability coordinates in the latter solve. In these
coordinates the maximum absolute row sum gives a Lipschitz bound.
A segment splits into finitely many intervals lying in the simplices
(intersect the segment with their rational affine half-spaces). Summing
their lengths along this straight line proves that the maximum bound is
valid globally. For a zero-dimensional source use zero.
Call the resulting bounds $A$ for $F$ and $B$ for $G$. Input in ambient
sup norm bounds input in the independent coordinates as well.

Set $M=AB$. On the nerve $z$ is $M$-Lipschitz, the numerator of $W_i$ is
$(1+M)$-Lipschitz, and its denominator is $m(1+M)$-Lipschitz. Since
$0\le y_i z_i\le1$ and the denominator is at least $1/m^2$, the identity
$a/b-a'/b'=(a-a')/b+a'(b'-b)/(bb')$ gives the bound
$C=(1+M)(m^2+m^5)$ for $W$. For the two-stage nerve homotopy use
$\max(1,C,M)+2$ on the product with time. At fixed time the interpolation
is a convex combination of the given maps, and at a fixed point its
speed is at most two. For the spatial homotopy use
$\max(1,M)+D$, with $D$ the largest coordinate span of the whole bounding
grid (zero when $d=0$). A bound $J$ gives input demand
$k+2+\lceil\log_2\max(1,J)\rceil$ for output scale $k$,
with room for rational rounding. These are closed under PR composition.
The finite subdivision searches, linear solves by rational determinants,
and maxima run over input-bounded tables. This proves the identities and
moduli first on rational tables with quantifier-free loop invariants. Continuity gives their regular-name interpretations by taking approximants
at the stated demand. No name is required to decide which side of a
grid face its limit occupies. The companion \texttt{box\_homotopy.py}
evaluates both maps and homotopies exactly on rational examples.

\paragraph{Finite filtrations.}
Suppose a finite sequence of such families satisfies
$B_{t,i}\subseteq B_{t+1,i}$ with fixed chart labels. Sort the endpoints
from all stages together. Inclusion now gives literal inclusions of the
three complexes $K_t,E_t,L_t$. Both projections commute with inclusion on
each basis cell. On homology their induced maps are isomorphisms, so
$(\pi_L)_*(\pi_K)_*^{-1}$ commutes with every stage map, by multiplying the two commuting squares by the inverses. Thus the inclusion maps of the unions
are also recovered from the nerves. The proof requires no coherent choice
of local contractions across stages. Every finite prefix admits this construction.

\paragraph{Checks and failure cases.}
The four sides of a rectangular frame thickened to overlapping strips give
the four-cycle in Figure~\ref{fig:frame}. The six closed square faces of a cube, used
as degenerate three-dimensional boxes, have a nerve with the eight triangles
of an octahedral boundary. The provided program checks $DH+HD=1$ on all
296 augmented basis columns for the frame and all 248 for the cube. Their
Betti vectors through degree two are $(1,1,0)$ and $(1,0,1)$.
For a counterexample to dropping the local test, cover a square boundary by
its two length-two arcs from vertex $0$ to opposite vertex $2$. Their
intersection consists of the two endpoints. Its augmented cycle $[0]+[2]$
has no edge filler. The nerve is an interval, whereas the union has a
one-dimensional homology class. The program returns the failed overlap.

\begin{figure}[htbp]
\centering
\begin{tikzpicture}[x=.82cm,y=.82cm,font=\small]
  \fill[accent!22] (0,0) rectangle (3,3);
  \fill[white] (1,1) rectangle (2,2);
  \foreach \x in {0,2} \foreach \y in {0,2}
    {\fill[accent!52] (\x,\y) rectangle ++(1,1);}
  \draw[ink] (0,0) rectangle (3,3);
  \draw[ink] (1,1) rectangle (2,2);
  \foreach \x in {1,2} {\draw[ink!60] (\x,0)--(\x,3);}
  \foreach \y in {1,2} {\draw[ink!60] (0,\y)--(3,\y);}
  \node at (1.5,.5) {1}; \node at (2.5,1.5) {2};
  \node at (1.5,2.5) {3}; \node at (.5,1.5) {4};
  \node[anchor=north] at (1.5,-.15) {Four closed strips};
  \begin{scope}[shift={(6.5,1.5)}]
    \coordinate (a) at (0,-1.15); \coordinate (b) at (1.15,0);
    \coordinate (c) at (0,1.15); \coordinate (d) at (-1.15,0);
    \draw[accent,thick] (a)--(b)--(c)--(d)--cycle;
    \foreach \pt in {a,b,c,d} {\fill[ink] (\pt) circle (2.5pt);}
    \node[below] at (a) {1}; \node[right] at (b) {2};
    \node[above] at (c) {3}; \node[left] at (d) {4};
    \node[anchor=north] at (0,-1.65) {Intersection nerve};
  \end{scope}
\end{tikzpicture}
\caption{Each darker square is an intersection of two strips and gives one edge of the nerve. There are no triple intersections. Both spaces have one independent loop.}
\label{fig:frame}
\end{figure}


\subsection{Arithmetic and Lean realization}

With $s$ cells and $m$ charts the incidence table has at most $s2^m$
ordinary columns. The added columns number at most $s$ or $2^m$.
A matrix can be packed as a list of bit masks of these bounded lengths.
For boxes, $s\le (4m)^d$ before removing cells outside the union.
Forming powers, testing local fillings, XOR summation, and sorting endpoints
by integer cross-products are primitive recursive in the full input,
including dimension. The outer correctness assertions quantify universally
over input codes. Inside a loop, assert by a Boolean test that all
previous columns obey their identity. For the power loop, test that all
nonzero entries lower the specified coordinate degree by at least the
current exponent. The indices in these tests range only over the two
explicit table sizes. In the telescoping identity one may compare each
matrix column, because equality of linear maps means equality of these
columns. Replacing bounded tests by their primitive recursive folds leaves
quantifier-free induction. A local test failure is an explicit alternative
of the total algorithm.

The Lean file \texttt{TopologyCase.lean} proves the nilpotent correction
identity for an abelian group of exponent two, including commutation of the
defect with $D$. Its bounded-lowering lemma derives nilpotence from a
strict decrease in the largest surviving opposite-coordinate degree. Its
common-coordinate bound proves the denominator estimate for $W$ after
clearing rational denominators. The executable \texttt{nerve.py} implements both local
constructions. To reach a full object derivation, represent a cell by a face
mask or a tuple of grid indices, implement the three arrays as bit matrices,
and instantiate in order the two-deletion cancellation, local-fill
induction, degree-lowering induction, and telescoping identity above.
Project the three block equations by bounded row and column indices before
composing the final maps. These are the remaining lemmas between the host
identity and the arithmetized chain equivalence.
